\section{Geometry of a signed boundary representative}
\label{sg:section}

The induction requires an abundance criterion for a nef log canonical
divisor whose restriction to the reduced boundary is semiample.
The representative supported on that boundary may have either sign.
We first state the result, then construct the geometry needed to lift
functions from the positive part of the boundary.

\begin{theorem}[Signed representative]\label{thm:signed}
Let $(X,D)$ be a projective $\Q$-factorial dlt pair over $\C$, with
$D=\sum_iD_i$ reduced.  Suppose that
\[
 L=K_X+D\quad\text{is nef},\qquad L-cD\quad\text{is pseudo-effective}
 \quad\text{for some }c>0,
\]
and that
\[
 L\sim_{\Q}G=\sum_i a_iD_i,\qquad a_i\in\Q.
\]
If $L|_D$ is semiample as a rational line bundle on the reduced scheme
$D$, then $L$ is abundant:
\[
 \kappa(X,L)=\nu(X,L).
\]
\end{theorem}

Put $n=\dim X$ and $v=\nu(X,L)$. In the nontrivial range $0<v<n$,
the geometric goal is a dominating family of $(n-v)$-dimensional
projective subvarieties disjoint from $D$. We will show that a suitable
positive-coefficient component has semiample image of dimension $v-1$.
Over a general point of that image, the construction in this section
separates the positive boundary from the negative and coefficient-zero
parts and makes it Cartier, with specified adjunction data.
Proposition~\ref{prop:formal-lifting} then lifts the image parameters
together with a transverse parameter through all infinitesimal
neighborhoods. Fixing the former and varying the latter will deform a
boundary fiber off $D$. Section~\ref{sg:descent} algebraizes these
deformations and descends $L$ along its nef reduction to prove abundance.
This argument uses no Iitaka subadditivity assumption.

\subsection{The positive boundary and its small intersections}

If $D=0$, the signed relation gives $L\sim_{\Q}0$, so the theorem
is immediate.  For the construction below we may therefore assume
$D\ne0$; its semiample restriction then defines the stated morphism
to a projective space.

Write
\[
 G=G_+-G_-,\qquad D_0=\sum_{a_i=0}D_i,
\]
where $G_+$ and $G_-$ are effective and have disjoint prime supports.
Take a positive integer $m$ such that $mG_\pm$, $mD_0$, and $mL$
are integral Cartier divisors and
\[
 N_0=\OO_X(mG)\simeq\OO_X(mL).
\]
After increasing $m$, we may assume that $N_0|_D$ is generated by
global sections.  Fix the morphism it defines,
\[
 \phi:D\longrightarrow P=\PP^s,\qquad
 N_0|_D\simeq\phi^*\OO_P(1),
\]
and write $s_0$ for the rational section of $N_0$ whose divisor is
$mG$. Fix an ample Cartier divisor $H$.
The next lemma locates the part of the boundary over
which the lifting construction will take place.

\begin{lemma}\label{sg:numerical-boundary}
If $v=n$, then $L$ is big.  If $v=0$, then $D=0$ and
$L\sim_{\Q}0$.  In the remaining case $0<v<n$, put $r=v-1$.
Every component of $D$ has $\phi$-image of dimension at most $r$,
and some component of $G_+$ has image of dimension $r$.
Moreover,
\[
 \dim\phi\bigl(\Supp G_+\cap(\Supp G_-\cup D_0)\bigr)<r;
\]
when $r=0$, the intersection in this formula is empty.
\end{lemma}

\begin{proof}
For $v=n$, the numerical criterion for a nef divisor gives bigness.
Assume $v<n$.  Intersect the pseudo-effective class $L-cD$ with the
indicated product of nef classes:
\[
 0\le (L-cD)L^vH^{n-v-1}
    =-c\sum_iD_iL^vH^{n-v-1}.
\]
Each number $D_iL^vH^{n-v-1}$ is nonnegative, so all of them vanish.
When $v=0$, ampleness then forces $D=0$; the signed representative
consequently gives $L\sim_{\Q}0$.

For $v>0$, semiampleness identifies the numerical dimension of the
restriction to each component of $D$ with that component's image
dimension under $\phi$.  The vanishing just proved bounds this
dimension by $r$.  Moreover,
\[
 0<L^vH^{n-v}=\sum_i a_iD_iL^rH^{n-v},
\]
so the bound is attained by a component with positive coefficient.

To control where that positive part meets the rest of the boundary,
consider the symmetric matrix
\[
 Q_{ij}=D_iD_jL^rH^{n-r-2}.
\]
Distinct effective $\Q$-Cartier divisors have effective intersection
cycles, so the off-diagonal entries are nonnegative.  The ambient
intersection form has at most one positive direction by the mixed
Hodge index theorem: approximate by ample classes and apply the
ordinary Hodge index theorem on complete-intersection surfaces.
In this form $L$ is isotropic and orthogonal to every $D_i$, while
its pairing with $H$ is strictly positive. Indeed, writing $Q$
also for the ambient bilinear intersection form, we have
\[
 Q(L,L)=L^{r+2}H^{n-r-2}=0,\qquad
 Q(L,D_i)=D_iL^{r+1}H^{n-r-2}=0,
\]
whereas $Q(L,H)=L^{r+1}H^{n-r-1}>0$ because $r+1=v$.
The form on the orthogonal space to $L$ is therefore negative
semidefinite. In particular this holds on the span of the $D_i$.
Pulling back
to a resolution gives the same conclusion, so smoothness of $X$
is unnecessary here.

Write $a=(a_i)$ for the coefficient vector. The signed relation
gives the vector equation $Qa=0$: for every $i$,
\[
 \sum_jQ_{ij}a_j=D_iL^{r+1}H^{n-r-2}=0.
\]
Temporarily write
$a=a^+-a^-$ for its decomposition into positive and negative
coefficient vectors.  We obtain
\[
 Q(a^+,a^+)=Q(a^+,a^-)\ge0.
\]
Negative semidefiniteness forces both sides to vanish. A vector
on which a negative semidefinite quadratic form vanishes belongs
to its kernel, so $Qa^+=0$. In a row with $a_i\le0$, the diagonal
term is absent and every summand $Q_{ij}a_j^+$ is nonnegative.
Their sum is zero; hence each intersection with a positive
component has zero $L^r$-degree against $H^{n-r-2}$. Semiampleness
on $D$ then makes its image dimension less than $r$.  For $r=0$,
any nonzero effective intersection would have positive ample
degree, so the intersection is empty.
\end{proof}

For the rest of the construction assume $0<v<n$, and put $N=n-1$.

\subsection{The divisor and adjunction data for lifting}

We will replace the positive boundary by a reduced Cartier divisor
whose normal line comes from the semiample image.  Its powers must
be actual line bundles throughout the construction.  For a positive
integer $\ell$, the root gerbe
\[
 \mathcal P=\sqrt[\ell]{\OO_P(1)}
\]
parametrizes $\ell$-th roots of the indicated line bundle without
the choice of a section.  Denote its tautological line by $C$;
thus $C^\ell$ is the pullback of $\OO_P(1)$.  Further pullbacks
of $C$ will carry the same notation.
The use of normalized cyclic covers and their eigenspace
decompositions follows the classical covering method of
Esnault--Viehweg
\cite[Section~3.5, Claim~3.10 and Corollary~3.11]
{EsnaultViehweg1992}.  Here we also retain the negative and
coefficient-zero boundary components and keep track of the fixed
adjunction isomorphism on the quotient charts.
We will also place the Stein image in a smooth target: this allows
differential operators and projective Hodge-module direct images to be
used on ordinary scheme charts. The infinitesimal neighborhoods will
later be cut down to a projective fiber before they are algebraized.

\begin{proposition}\label{prop:signed-construction}
Choose $\ell$ divisible by $m$ and every nonzero integer $|ma_i|$,
and set $a=\ell/m$.  There is a normal tame Deligne--Mumford stack
$\mathcal X$, considered on a neighborhood of a reduced Cartier
divisor $S$ of pure dimension $N=n-1$, together with a morphism
to $X$ and a reduced boundary
$T$ having no component in common with $S$, with the following
properties.
\begin{enumerate}[label=\textup{(\roman*)}]
\item
The pair $(\mathcal X,S+T)$ is log canonical, and a fixed
isomorphism of reflexive sheaves is given by
\[
 \omega_{\mathcal X}(S+T)\simeq\OO_{\mathcal X}(aS).
 \tag{\(\ast_{\mathrm{adj}}\)}\label{sg:ambient-adjunction}
\]
The support of $T$ is locally set-theoretically principal.
The ambient space and $S$ are Cohen--Macaulay on scheme charts.
Off $T$, the ambient space is Gorenstein and the displayed
isomorphism is ordinary Cartier adjunction data.

\item
Put $J=(S\cap T)_{\mathrm{red}}$.  Both $S$ and $J$ are Du Bois
on scheme charts, the ideal $\mathcal I_{J/S}$ is maximal
Cohen--Macaulay, and
\[
 \mathcal B:=
 \mathcal Hom_S(\mathcal I_{J/S},\omega_S)
 \simeq\OO_S(aS).
\]
The identification agrees off $J$ with the residue convention
specified by \eqref{sg:ambient-adjunction}.

\item
There is a finite representable morphism
\[
 S\longrightarrow D\times_P\mathcal P
\]
whose image covers $\Supp G_+$.  The line $\OO_S(S)$ is the pullback
of $C$.  The image of $J$ in $P$ has dimension less than $r$,
whereas the image of $S$ has dimension $r$.

\item
There is a smooth projective bundle $p:\mathcal Y\to\mathcal P$
and a representable projective morphism $g:S\to\mathcal Y$.
This is a morphism from the reduced divisor $S$; no extension
to a neighborhood of $S$ in $\mathcal X$ is asserted here.
Writing $h=p\circ g$ and
$\mathcal T=\Spec_{\mathcal P}h_*\OO_S$, the morphism $g$ factors
through a closed embedding $\mathcal T\hookrightarrow\mathcal Y$,
and
\[
 g_*\OO_S=\OO_{\mathcal T},\qquad
 \OO_S(S)=g^*C,\qquad \mathcal B=g^*C^a.
\]

\item
For a separated quasi-projective scheme chart
$U\to\mathcal Y$ that is etale and of finite type, put
$S_U=S\times_{\mathcal Y}U$.
The etale morphism $S_U\to S$ extends compatibly to every
nilpotent thickening $qS$ in $\mathcal X$.
Only this etale chart is extended; an extension of $g$ to $qS$
is not needed to define the chart.
The resulting $qS_U$ are quasi-projective schemes, separated
and quasi-finite over $X$.
\end{enumerate}
All the adjunction and duality identifications are compatible
with changes of etale chart.
\end{proposition}

\begin{proof}
We construct the divisor and its adjunction data first, then place
its Stein image in a smooth projective bundle.  The final step
verifies that the infinitesimal neighborhoods have the scheme
charts required for formal lifting.

\smallskip\noindent\emph{Roots and a fixed adjunction isomorphism.}
Let $\rho:\mathcal R\to X$ be the normalized simultaneous root stack of the Cartier
divisors $mG_+$, $mG_-$, and $mD_0$, taking an $\ell$-th root of
each divisor together with its section; the construction allows
empty divisors.  Denote the tautological root divisors by
$S_+,S_-,S_0$ on $\mathcal R$. Their reducedness can be checked
at a generic prime of downstairs multiplicity $b$. The chosen
divisibility conditions give $b\mid\ell$. A normalized chart of
$y^\ell=x^b$ has ramification index $\ell/b$, and $y$ has order
one.  Since $m$ divides $\ell$, this also applies to $D_0$.
The divisors are Cartier, so normality and generic reducedness
give their reducedness everywhere.

Put $B=S_++S_-+S_0$ on $\mathcal R$. With this full reduced
boundary, log ramification gives a log
canonical pair and the relation
\[
 K_{\mathcal R}+B\sim_{\Q}a(S_+-S_-).
\]
The ambient charts are klt.  They are klt off the boundary,
and, on a log resolution, decreasing every boundary coefficient
slightly removes all zero-discrepancy places. The difference
between the two sides is a torsion integral Weil class, represented by
\[
 \mathcal F=
 \bigl(\omega_{\mathcal R}(B)
       \otimes\OO_{\mathcal R}(-a(S_+-S_-))\bigr)^{**}.
\]
Choose a periodicity isomorphism $\mathcal F^{[q]}\simeq\OO_{\mathcal R}$,
where square brackets denote reflexive powers. It defines multiplication
in the finite algebra
\[
 \mathcal A=\bigoplus_{i=0}^{q-1}\mathcal F^{[i]}.
\]
Let $\tau:\mathcal R^\dagger\to\mathcal R$ be the normalization
of $\Spec_{\mathcal R}\mathcal A$. In codimension one this is
the ordinary cover of a torsion line bundle and is etale. The
cover preserves log canonicity, klt ambient singularities, and
the reduced Cartier boundary. From now on $S_+,S_-,S_0$ denote
their pullbacks to $\mathcal R^\dagger$.

Tautological evaluation on this cover trivializes the pullback
of $\mathcal F$ in codimension one and hence fixes the adjoint
isomorphism. Because this evaluation is a sheaf morphism
on the cover stack itself, the adjoint isomorphism is equivariant
on every atlas.  This equivariance will be needed for the
quotient below.

\smallskip\noindent\emph{Separating the positive and negative parts.}
Let $p_1:\widetilde{\mathcal R}\to\mathcal R^\dagger$ be the
normalized blowup of the ideal of $S_+\cap S_-$. The spaces
constructed so far are related by
\[
 \widetilde{\mathcal R}\xrightarrow{p_1}\mathcal R^\dagger
 \xrightarrow{\tau}\mathcal R\xrightarrow{\rho}X.
\]
The blown-up ideal is
locally generated by two elements, so the ordinary blowup embeds
in a relative projective line.  Its fibers have dimension at
most one, as do the fibers after finite normalization.  Hence
each exceptional divisor lies over a codimension-two component
of the intersection.  Downstairs such a component is a dlt
stratum, generically SNC.  Separate normalized Kummer charts
and purity for the index cover at the smooth generic locus
give the same description upstairs.

The tautological exceptional divisor $E$ is consequently reduced
Cartier, and the divisors
\[
 S'_+=p_1^*S_+-E,\qquad S'_-=p_1^*S_--E
\]
are reduced Cartier and disjoint on $\widetilde{\mathcal R}$.
Since no codimension-two two-branch stratum is contained
in $S_0$, its pullback $p_1^*S_0$ is reduced Cartier and has no
exceptional component.  The total boundary
\[
 S'_++S'_-+E+p_1^*S_0
\]
is crepant and log canonical.  To check the ambient singularities
as well, denote this boundary on a chart $V'$ by $B'$.  The
crepant equality gives log canonicity for every valuation.
Outside $B'$, the blowup is an isomorphism to the old klt
boundary complement.  A divisor $F$ with
$a(F;V',B')=0$ therefore has center in $\Supp B'$.  Since $B'$
is effective Cartier, $\ord_F(B')>0$, whence
\[
 a(F;V',0)=a(F;V',B')+\ord_F(B')>0.
\]
Positive pair discrepancies also remain positive after the
boundary is dropped.  Thus $V'$ is klt even at its
higher-codimension singular loci. On $\widetilde{\mathcal R}$,
put $S=S'_+$ and $T=S'_-+E+p_1^*S_0$. Near $S$, the section of the disjoint
divisor $S'_-$ is a unit, and the fixed linear equivalence reads
\[
 \omega(S+T)\simeq\OO(aS).
\]
Klt singularities make the ambient charts Cohen--Macaulay.
All the boundary divisors in this display are Cartier, so the
isomorphism also makes these charts Gorenstein.

We also need finiteness of $S\to S_+$ on $\mathcal R^\dagger$.
Locally write $S_+=(u=0)$ and $S_-=(v=0)$. Before normalization,
the strict transform of $S_+$ lies in the blowup chart with
ratio $u/v$, where its equation is $u/v=0$. Thus each fiber of
this strict transform over $S_+$ has at most one point.
Properness and finite normalization prove finiteness of $S\to S_+$.
A codimension-one point of $S\cap T$ therefore lies over a
codimension-two intersection downstairs.  The generic SNC
description shows that $T|_S$ is generically reduced.  Being
Cartier on the Cohen--Macaulay scheme $S$, it is reduced
everywhere.

\smallskip\noindent\emph{Retaining the required root line.}
The construction so far has supplied reduced Cartier divisors
and fixed adjunction data.  We now remove the root characters
that are unnecessary, while preserving the line attached to $S$.
On $\widetilde{\mathcal R}$ the line $\OO(S-S'_-)$ has $\ell$-th power equal to the pullback
of $N_0$ and hence defines a morphism to the gerbe of $\ell$-th
roots of $N_0$ on $X$. Denote this gerbe by $\mathcal G$, and let
$\kappa:\widetilde{\mathcal R}\to\mathcal X$ be the relative
coarse-space morphism over $\mathcal G$. On a trivializing chart, we quotient
by the kernel of the finite-group character on the indicated
root line.  The quotients are ordinary quasi-projective schemes:
the covers used above are finite, the blowup is projective,
and finite quotients preserve quasi-projectivity.  These
chartwise quotients patch.

Before taking the quotient, the section of the disjoint divisor
$S'_-$ is a unit near $S$. There the retained line
$\OO(S-S'_-)$ identifies with $\OO(S)$, so the line of $S$ and
its section both descend through the kernel quotient. Continue
to write $S,T$ for their reduced images under $\kappa$; they
now lie on $\mathcal X$. In particular $S$
remains Cartier.  The finite-group norm of a local equation
for $T$ upstairs makes its image set-theoretically principal;
Cartierness of the reduced image $T$ is not needed.
Divisorial ramification occurs only on the boundary.  Log
ramification thus preserves log canonicity and identifies the
descended adjoint isomorphism with
\eqref{sg:ambient-adjunction}.  The ambient chart and $S$ are
Cohen--Macaulay because their structure sheaves are invariant
direct summands of the finite CM covers.

More explicitly, the kernel acts trivially on the retained root
line and hence on $\OO(aS)$.  Equivariance of the fixed
isomorphism gives the same kernel character on $\omega(S+T)$.
Taking invariants descends the isomorphism, and log ramification
identifies the invariant log dualizing sheaf.  Off $T$, this
isomorphism makes the canonical sheaf invertible, giving the
asserted Gorenstein property.

\smallskip\noindent\emph{Duality on the boundary.}
The reduced supports $S$ and $J$ are unions of log canonical
centers, since intersections of lc centers are again unions
of lc centers.  Both are therefore Du Bois by
\cite[Theorems~1.4 and~1.7]{KollarKovacs2010}.
To identify the ideal and its dual, work on a quotient chart
and denote the finite cover above by $f:S^{\mathrm{up}}\to S$.
Reducedness of $T^{\mathrm{up}}|_{S^{\mathrm{up}}}$ gives
\[
 \mathcal I_{J/S}
 =
 \bigl(f_*\OO_{S^{\mathrm{up}}}
       (-T^{\mathrm{up}}|_{S^{\mathrm{up}}})\bigr)^{\mathrm{inv}}.
\]
An invariant function vanishes on the reduced quotient image
exactly when its pullback vanishes on this reduced preimage.
The displayed invariant sheaf is maximal Cohen--Macaulay:
the upstairs sheaf is invertible on a CM scheme, finite
pushforward preserves its depth over the quotient, and in
characteristic zero invariants form a direct summand.

Apply finite-map duality with trace and then take invariants:
\[
 \mathcal Hom_S(\mathcal I_{J/S},\omega_S)
 =
 \bigl(f_*\omega_{S^{\mathrm{up}}}
       (T^{\mathrm{up}}|_{S^{\mathrm{up}}})\bigr)^{\mathrm{inv}}
 \simeq\OO_S(aS).
\]
The last isomorphism comes from Cartier adjunction upstairs
and descent of the line of $S$.  Finite-map duality applies
in the presence of ramification; it requires no invertibility
of $\omega_S$ itself.  Off $J$, the identification is the
ordinary residue fixed by the ambient isomorphism.  Normalized
trace preserves this residue convention, so the identifications
agree on overlaps.

\smallskip\noindent\emph{The target of the lifting argument.}
On $S$, the root gerbe of $N_0$ is $D\times_P\mathcal P$.
Finiteness of the strict divisor, proved above, and removal
of the relative inertia kernel give a finite representable
map $S\to D\times_P\mathcal P$.  Its image covers the positive
boundary, and the root line is $\OO_S(S)=C$.
Lemma~\ref{sg:numerical-boundary} therefore gives image dimension
$r$ for $S$ in $P$. To check the smaller image of $J$, recall
that before the quotient $S'_-$ is disjoint from $S$, the
exceptional divisor $E$ maps into $\Supp G_+\cap\Supp G_-$,
and $p_1^*S_0$ maps into $D_0$. Taking the quotient changes none
of these images in $X$. The image of $J$ in $D$ is consequently
contained in
\[
 \Supp G_+\cap(\Supp G_-\cup D_0).
\]
The same lemma gives image dimension less than $r$ for this
locus in $P$. We have thus obtained a
representable projective morphism $h:S\to\mathcal P$; it remains
to embed its Stein image in the required smooth bundle.

The Stein algebra defines a finite stack
$\mathcal T=\Spec_{\mathcal P}h_*\OO_S$ over $\mathcal P$.
This coherent algebra is generated as a module by vector
bundles on $\mathcal P$.  Indeed, decompose it by the finitely
many inertia characters, twist each summand by a power of $C$
so that it descends to $P$, and use Serre generation there.
The resulting vector-bundle surjection embeds $\mathcal T$ in
a vector bundle over $\mathcal P$.  Since $\mathcal T$ is
proper over the base, it remains closed in the projective
completion. This completion is $\mathcal Y$ and gives the map
$g:S\to\mathcal Y$ from the reduced divisor.
The Stein construction yields $g_*\OO_S=\OO_{\mathcal T}$,
and the line identities follow from those already established.

Finally, start with the etale chart $S_U\to S$ obtained from $g$.
Invariance of the etale site under nilpotent thickening extends
this chart uniquely to $qS_U\to qS$, compatibly as $q$ varies.
This construction does not use a map $qS\to\mathcal Y$;
lifting the map from $S$ is the task of the next section.
The reduction $S_U$ is a scheme; its extension $qS_U$ is an
algebraic space with trivial inertia.  The chosen chart is
separated, the map to the root gerbe is finite, and that gerbe
has finite diagonal, so $qS_U$ is separated over $X$.
The finite representable map $S\to D\times_P\mathcal P$
and the etale map $S_U\to S$ show that the finite-type
geometric fibers of $S_U\to X$ are zero-dimensional.
Nilpotent thickening changes neither geometric points nor
fiber dimensions. The morphism $qS_U\to X$
is therefore quasi-finite as well.  Zariski's main theorem
embeds $qS_U$ in a scheme finite over $X$, making it a
quasi-projective scheme as required.
\end{proof}
